\begin{examplebox}
\textbf{例 1：原题不变，选纵坐标计算更短。}
$y^2=4x$ 与 $x+2y-2=0$ 交于 $P,Q$，求弦长。

\textbf{解。} $2p=4$，所以 $p=2$；$A=1,B=2,C=-2$。
从抛物线得到 $x=y^2/4$，代入直线：
\[
 \frac{y^2}{4}+2y-2=0
 \ \Longleftrightarrow\ y^2+8y-8=0.
\]
因此 $y_1+y_2=-8$、$y_1y_2=-8$，
\[
 |y_2-y_1|=\sqrt{64+32}=4\sqrt6.
\]
直线 $x=2-2y$ 给出 $|x_2-x_1|=2|y_2-y_1|$，所以
\[
 PQ=\sqrt5\,|y_2-y_1|=\boxed{4\sqrt{30}}.
\]
\textbf{回验一：通用式。} $H=2^2\cdot2^2-2\cdot2\cdot1\cdot(-2)=24>0$，
$PQ=2\sqrt{5\cdot24}=4\sqrt{30}$。

\textbf{回验二：原横坐标固定形式。} 得到 $x^2-20x+4=0$，横坐标根差为 $\sqrt{384}=8\sqrt6$。
乘以 $\sqrt{1+(-1/2)^2}$，仍得相同弦长。
两个方向均可用，选系数更小、展开更短的方向即可。
\end{examplebox}
\clearpage
\begin{examplebox}
\textbf{例 2：向上开口，求弦中点。}
$x^2=8y$ 与 $2x-y+1=0$ 交于 $P,Q$，求中点 $S$。
\begin{center}\begin{tikzpicture}[font=\small]
\draw[black!30,rounded corners=2pt] (0,0) rectangle (14.4,8.2);
\node[anchor=west,font=\small\bfseries] at (.5,7.72) {中点先取横坐标平均，再代入弦所在直线};
% CURVE vertical 4 1
\begin{scope}[shift={(6,1.65)},scale=0.135]
\begin{scope}
\clip (-18,-2) rectangle (18,37);
\draw[blue,very thick,domain=-18:18,samples=181,variable=\t] plot (\t,{\t*\t/(8)});
\end{scope}
\draw[-Stealth,black!45] (-18,0)--(18,0) node[right] {$x$};
\draw[-Stealth,black!45] (0,-2)--(0,37) node[above] {$y$};
\draw[orange!85!black,very thick] (-1.5,-2)--(18,37);
\coordinate (P) at (-0.4852813742,0.0294372515);\node[fill=black,circle,inner sep=2pt] at (P) {};\node[below left=9pt] at (P) {$P$};
\coordinate (Q) at (16.4852813742,33.9705627485);\node[fill=black,circle,inner sep=2pt] at (Q) {};\node[left=12pt] at (Q) {$Q$};
\coordinate (S) at (8.0000000000,17.0000000000);\node[fill=violet,circle,inner sep=2pt] at (S) {};\node[left=12pt,text=violet] at (S) {$S$};
\end{scope}\node[anchor=west] at (.8,6.7) {$x^2=8y$};\node[anchor=west] at (.8,5.7) {$y=2x+1$};\node[anchor=west,violet] at (9.8,4.1) {$S=(8,17)$};
\node[anchor=west] at (.7,.55) {曲线与坐标等比例绘制；中点 $S$ 在 $PQ$ 上，但 $8^2\ne8\times17$};
\end{tikzpicture}\end{center}

\textbf{解。} $p=4$，由 $y=2x+1$ 得 $x^2=8(2x+1)$，即
\[
 x^2-16x-8=0.
\]
判别式 $288>0$，确有两点。韦达给出 $x_S=(x_1+x_2)/2=8$，
所以 $y_S=2\cdot8+1=17$，答案为 $\boxed{S=(8,17)}$。

\textbf{回验。} 交点为 $(8\pm6\sqrt2,\,17\pm12\sqrt2)$，同号配对，分别满足两条原方程；逐坐标取平均即为 $(8,17)$。
中点在直线上，不要求也在抛物线上。
\end{examplebox}
\clearpage
\begin{examplebox}
\textbf{例 3：中点条件的候选值，必须用两点条件筛选。}
$y^2=4x$ 与 $y=kx+1$ 交于两个不同点，且弦中点在 $x=1$ 上，求 $k$。
\begin{center}\begin{tikzpicture}[font=\small]
\draw[black!30,rounded corners=2pt] (0,0) rectangle (14.4,8.2);
\node[anchor=west,font=\small\bfseries] at (.5,7.72) {中点方程给出候选值，仍须排除相切};
% CURVE lateral 2 1
\begin{scope}[shift={(5.8,3.9)},scale=0.8]
\begin{scope}
\clip (-0.5,-3.1) rectangle (3,3.1);
\draw[blue,very thick,domain=-3.1:3.1,samples=181,variable=\t] plot ({\t*\t/(4)},\t);
\end{scope}
\draw[-Stealth,black!45] (-0.5,0)--(3,0) node[right] {$x$};
\draw[-Stealth,black!45] (0,-3.1)--(0,3.1) node[above] {$y$};
\draw[violet,very thick] (-0.5,0.5)--(2.1,3.1);
\draw[orange!85!black,very thick] (-0.5,2)--(2.05,-3.1);
\draw[teal!70!black,thick,densely dashed] (1,-3.1)--(1,3.1);
\coordinate (T) at (1.0000000000,2.0000000000);\fill[violet] (T) circle (2pt);
\coordinate (S) at (1.0000000000,-1.0000000000);\fill[violet] (S) circle (2pt) node[right=11pt] {$S$};
\coordinate (P) at (0.1339745962,0.7320508076);\fill[black] (P) circle (2pt) node[above left=10pt] {$P$};
\coordinate (Q) at (1.8660254038,-2.7320508076);\fill[black] (Q) circle (2pt) node[above right=10pt] {$Q$};
\draw[-Stealth,black!45,thin] (1.3,3.4)--(T);\node[above right=3pt] at (1.3,3.4) {$T$};\end{scope}\node[anchor=west,violet] at (9.5,6.5) {$k=1$：相切};\node[anchor=west,orange!85!black] at (9.5,5.5) {$k=-2$：两点};\node[anchor=west,teal!70!black] at (9.5,3.8) {$x_S=x_T=1$};
\node[anchor=west] at (.7,.55) {$y^2=4x,\ y=kx+1$；$T$ 是重合端点，真正非零弦的中点是 $S$};
\end{tikzpicture}\end{center}

\textbf{解（中点短法）。} $k=0$ 时只有一个交点，不合题意。
$k\ne0$ 时，$A=k,B=-1,C=1,p=2$，由已证公式 $y_S=2/k$。
又 $S$ 在直线且 $x_S=1$，故 $2/k=k+1$，得到 $k=1$ 或 $-2$。
真正两点还需 $H=4-4k>0$，即 $k<1$，所以 $\boxed{k=-2}$。

\textbf{回验。} $k=1$ 时 $(x+1)^2=4x$，只有重根 $x=1$，是切点 $T(1,2)$。
$k=-2$ 时得到 $4x^2-8x+1=0$，判别式 $48>0$；横坐标和为 $2$，中点恰为 $S(1,-1)$，符合题意。
\end{examplebox}
\clearpage
